\chapter{Losowe miary skoków i całkowanie względem miar Poissona}
\label{ch:poisson-random-measure}
\paragraph{Odsyłacze międzytomowe.} Proces Poissona z tomem II (rozdział~\ref{II-ch:poisson-process}) łączy się ze statystyczną estymacją liczby zdarzeń z tomu IV (rozdział~\ref{IV-ch:conjugate-families}).


\section{Dlaczego nie wystarcza suma skoków?}

Złożony proces Poissona zapisaliśmy jako $J_t=\sum_{j=1}^{N_t}Y_j$.
W tym wzorze ukryte są \emph{dwie} informacje: czas wystąpienia
skoku i jego rozmiar. Gdy skoków jest nieskończenie wiele, sumowanie
wszystkich ich rozmiarów może nie mieć sensu, nawet na przedziale
skończonym. Warunek definiujący miarę Lévy'ego gwarantuje tylko
całkowalność \emph{kwadratów} małych skoków. Dlatego rozdzielamy
losowe zliczanie skoków od ich całkowania, a następnie odejmujemy
od części sum odpowiednie przewidywalne średnie.

Niech $E$ będzie przestrzenią polską z borelowskim $\sigma$-ciałem
$\calB(E)$, a $\nu$ miarą $\sigma$-skończoną na $E$.
W przykładach głównych $E=\R^d\setminus\{0\}$ i $\nu$ będzie
miarą Lévy'ego. Pracujemy na ustalonym horyzoncie $T>0$;
$(\calF_t)_{0\le t\le T}$ oznacza filtrację spełniającą warunki
zwykłe. Używamy zapisu $\dd t\,\dd\nu(z)$ dla miary
produktowej $\lambda_1\otimes\nu$.

\begin{definicja}[Losowa miara Poissona]
\label{def:PRM}
Losową miarą Poissona o intensywności $\dd t\,\dd\nu(z)$ na
$(0,T]\times E$ nazywamy odwzorowanie $N$, które dla prawie każdego
$\omega$ jest miarą zliczającą i spełnia:
\begin{enumerate}
\item dla każdego mierzalnego $B\subset(0,T]\times E$ o
$(\lambda_1\otimes\nu)(B)<\infty$ zmienna $N(B)$ ma rozkład
Poissona o parametrze $(\lambda_1\otimes\nu)(B)$;
\item dla każdej skończonej rodziny parami rozłącznych takich
zbiorów $B_1,\ldots,B_m$ zmienne $N(B_1),\ldots,N(B_m)$ są niezależne;
\item $N((0,t]\times A)$ jest $\calF_t$-mierzalna, a
$N((s,t]\times A)$ jest niezależna od $\calF_s$ dla $s<t$
i $\nu(A)<\infty$.
\end{enumerate}
Przez $\widetilde N(\dd t,\dd z)$ oznaczamy formalnie
$N(\dd t,\dd z)-\dd t\,\dd\nu(z)$. Sens całkowy tej różnicy
określimy niżej.
\end{definicja}

\begin{uwaga}[O zbiorach nieskończonej intensywności]
Jeżeli $\nu(A)=\infty$, to $N((0,t]\times A)=\infty$
prawie na pewno dla $t>0$, o ile taka miara Poissona istnieje.
Nie oznacza to, że $\widetilde N((0,t]\times A)$ jest określoną
różnicą dwóch nieskończoności: \emph{nie definiujemy} jej w ten
sposób. Kompensowaną miarę rozumiemy przez całki po zbiorach o
skończonej intensywności i przedłużenie w $L^2$.
\end{uwaga}

\begin{twierdzenie}[Istnienie i konstrukcja przez znaczniki]
\label{th:prm-existence}
Dla każdej $\sigma$-skończonej miary $\nu$ na $E$ istnieje
losowa miara Poissona na $(0,\infty)\times E$ o intensywności
$\dd t\,\dd\nu(z)$, z naturalną uzupełnioną filtracją
zapewniającą warunek niezależności przyrostów z
Definicji~\ref{def:PRM}.
\end{twierdzenie}
\begin{proof}
\emph{Krok 1: skończona miara.} Niech $0<\nu(E)=\lambda<\infty$.
Bierzemy proces Poissona $K$ o intensywności $\lambda$, jego
chwile skoków $S_j$ oraz niezależne znaczniki $Z_j\in E$ o
wspólnym prawie $\rho=\nu/\lambda$, niezależne od $K$.
Definiujemy $N=\sum_{j\ge1}\delta_{(S_j,Z_j)}$.
Dla rozłącznych $A_1,\ldots,A_m\subset E$ i przedziału
$(s,t]$ warunkowa funkcja tworząca wektora
$(N((s,t]\times A_k))_k$ wynosi
$(\sum_{k=1}^m\rho(A_k)u_k+1-\sum_{k=1}^m\rho(A_k))^{K_t-K_s}$.
Po uśrednieniu po rozkładzie Poissona dostajemy
$\exp((t-s)\sum_k\nu(A_k)(u_k-1))$.
Jest to iloczyn funkcji tworzących niezależnych zmiennych
Poissona o parametrach $(t-s)\nu(A_k)$.
Dalsze rozbicie przedziału czasu pokazuje niezależność na
rozłącznych prostokątach czasowo-przestrzennych; z twierdzenia
o klasie monotonicznej rozciągamy to na rozłączne zbiory
mierzalne o skończonej intensywności. Obliczenia pokazują
również niezależność od informacji do chwili $s$.

\emph{Krok 2: miara $\sigma$-skończona.} Wybierzmy rozkład
$E=\bigsqcup_{k\ge1}E_k$ z $\nu(E_k)<\infty$.
Na iloczynie przestrzeni probabilistycznych konstruujemy niezależne
miary $N_k$ o intensywnościach $\dd t\,\nu|_{E_k}$ i kładziemy
$N(B)=\sum_kN_k(B\cap((0,\infty)\times E_k))$.
Dla $(\lambda_1\otimes\nu)(B)<\infty$ suma parametrów Poissona
jest skończona, więc $N(B)<\infty$ prawie na pewno.
Niezależność dla rozłącznych zbiorów przechodzi z sum skończonych
przez zbieżność funkcji tworzących, a adaptację zapewnia naturalna
filtracja wygenerowana przez wszystkie zliczenia do danej chwili.
Uzupełnienie i prawostronna regularizacja tej filtracji zachowują
niezależność przyrostów: dla chwil $s_n\downarrow s$ niezależność
od $\calF^0_{s_n}$ przechodzi na $\cap_n\calF^0_{s_n}$, ponieważ
zliczenia na $(s,s_n]$ dążą do zera według prawdopodobieństwa
na zbiorach o skończonej intensywności.
\end{proof}

\section{Średnia liczby skoków i formuła kompensacji}

\begin{twierdzenie}[Formuła Campbella i kompensacji]
\label{th:campbell}
Niech $N$ będzie losową miarą Poissona z Definicji
\ref{def:PRM}. Dla każdej nieujemnej, mierzalnej funkcji
$g:(0,T]\times E\to[0,\infty]$ mamy
\[
\E\int_{(0,T]\times E}g(t,z)\,N(\dd t,\dd z)
=\int_0^T\int_Eg(t,z)\,\dd\nu(z)\,\dd t.
\]
Jeżeli $H(\omega,t,z)\ge0$ jest przewidywalna względem
$(\calF_t)$ w zmiennej czasowej i borelowska w $z$, to
\[
\E\int H\,\dd N
=\E\int_0^T\int_E H_t(z)\,\dd\nu(z)\,\dd t.
\tag{11.1}\label{eq:compensation}
\]
Równość dopuszcza wartość $+\infty$ po obu stronach.
\end{twierdzenie}
\begin{proof}
Najpierw $g=\ind{(s,t]\times A}$, gdzie $\nu(A)<\infty$:
obie strony wynoszą $(t-s)\nu(A)$. Liniowość daje
wzór dla nieujemnych funkcji prostych, a twierdzenie o
zbieżności monotonicznej dla dowolnych $g$.

Dla losowego $H$ zaczynamy od
$H_r(z)=\xi\ind{(s,t]}(r)\ind A(z)$, gdzie $\xi$ jest
nieujemna i $\calF_s$-mierzalna, a $\nu(A)<\infty$.
Niezależność $N((s,t]\times A)$ od $\calF_s$ implikuje
$\E[\xi N((s,t]\times A)]=\E\xi\,(t-s)\nu(A)$.
Liniowość daje tezę dla sum takich procesów. Ich dodatnie
kombinacje oraz granice monotoniczne generują wszystkie
nieujemne funkcje przewidywalne na
$\Omega\times(0,T]\times E$, ponieważ prostokąty
$F_s\times(s,t]\times A$, $F_s\in\calF_s$,
generują przewidywalne $\sigma$-ciało (po domknięciu
względem prostych funkcji przestrzennych). Zatem
klasa monotoniczna i Tonelli kończą dowód.
\end{proof}

\begin{przyklad}[Ile jest dużych skoków?]
Dla miary Lévy'ego $\nu$ i $\varepsilon>0$ liczba skoków
rozmiaru większego od $\varepsilon$ w czasie $t$ jest
rozkład Poissona o parametrze $t\nu(\{\|z\|>\varepsilon\})$.
Prawdopodobieństwo braku takiego skoku wynosi
$\exp[-t\nu(\{\|z\|>\varepsilon\})]$.
Jeżeli $\nu(\{0<\|z\|\le1\})=\infty$, to dla każdego
$t>0$ występuje nieskończenie wiele małych skoków
prawie na pewno, choć suma ich kwadratów ma skończoną
wartość oczekiwaną, gdyż
$\E\sum_{s\le t}\|\Delta X_s\|^2\ind{\{\|\Delta X_s\|\le1\}}
=t\int_{\|z\|\le1}\|z\|^2\,\dd\nu(z)$ w konstrukcji
procesu Lévy'ego.
\end{przyklad}

\section[Całka skompensowana: konstrukcja w L2]{Całka skompensowana: konstrukcja w przestrzeni $L^2$}

\begin{definicja}[Elementarny proces przewidywalny]
\label{def:prm-step}
Proces $H(\omega,t,z)$ na $[0,T]\times E$ jest
\emph{elementarny}, jeżeli ma postać skończonej sumy
\[
H_t(z)=\sum_{j=0}^{m-1}\sum_{k=1}^{r}
\xi_{jk}\ind{(t_j,t_{j+1}]}(t)\ind{A_k}(z),
\]
gdzie $0=t_0<\cdots<t_m=T$, zbiory $A_k\in\calB(E)$
są parami rozłączne i mają skończoną miarę $\nu$,
a $\xi_{jk}$ są ograniczonymi, $\calF_{t_j}$-mierzalnymi
zmiennymi losowymi. Dla takiego procesu definiujemy
$\int_0^t\int_EH_s(z)\,\widetilde N(\dd s,\dd z)$
przez zastąpienie każdego zliczenia $N((t_j,t_{j+1}]\times A_k)$
jego wartością po odjęciu $(t_{j+1}-t_j)\nu(A_k)$;
na odcinku końcowym stosujemy $t\wedge t_{j+1}$.
\end{definicja}

\begin{twierdzenie}[Izometria skompensowanej całki Poissona]
\label{th:prm-isometry}
Niech $H,K$ będą elementarnymi procesami przewidywalnymi.
Jeżeli $I_t(H)=\int_0^t\int_EH_s(z)\,\widetilde N(\dd s,\dd z)$,
to $I(H)$ jest kwadratowo całkowalnym martyngałem o
prawostronnie ciągłych trajektoriach i
\[
\E I_T(H)I_T(K)
=\E\int_0^T\int_E H_t(z)K_t(z)\,\dd\nu(z)\,\dd t.
\tag{11.2}\label{eq:prm-isometry}
\]
W szczególności $\E I_T(H)^2=\E\int H^2\,\dd t\,\dd\nu$.
\end{twierdzenie}
\begin{proof}
Wspólnie rozdrobnijmy podziały czasu i rozbijmy zbiory
przestrzenne na rozłączne atomy. Oznaczmy
$\Delta\widetilde N_{jk}=N((t_j,t_{j+1}]\times A_k)
-(t_{j+1}-t_j)\nu(A_k)$.
Zmienne te mają średnią zero, wariancję
$(t_{j+1}-t_j)\nu(A_k)$; zmienne odpowiadające
rozłącznym atomom są niezależne. Dla $j<l$ wyraz
$\xi_{jk}\eta_{lq}\Delta\widetilde N_{jk}
\Delta\widetilde N_{lq}$ ma wartość oczekiwaną zero:
wszystkie czynniki poza ostatnim są $\calF_{t_l}$-mierzalne,
a ostatni ma warunkową średnią zero. Dla tego samego $j$
i różnych $k,q$ znika warunkowa kowariancja ze względu
na $\calF_{t_j}$. Zostają wyłącznie wyrazy diagonalne,
z których otrzymujemy
$\E[\xi_{jk}\eta_{jk}(t_{j+1}-t_j)\nu(A_k)]$.
Po zsumowaniu daje to wzór \eqref{eq:prm-isometry}.
Z tej samej warunkowej średniej zero wynika martyngałowość
względem filtracji. Na każdym przedziale elementarny
proces całkowy jest skończoną sumą kompensowanych
procesów zliczających, więc ma trajektorie càdlàg.
\end{proof}

\begin{twierdzenie}[Rozszerzenie na $L^2$ procesów przewidywalnych]
\label{th:prm-extend}
Niech $H$ będzie procesem przewidywalnym i
$\E\int_0^T\int_E |H_t(z)|^2\,\dd\nu(z)\,\dd t<\infty$.
Istnieje jednoznaczny (z dokładnością do nierozróżnialności)
càdlàg martyngał $I(H)$, będący granicą ucp całek
procesów elementarnych $H^{(n)}$ takich, że
$\E\int|H^{(n)}-H|^2\,\dd t\,\dd\nu\to0$.
Zachodzi izometria \eqref{eq:prm-isometry} dla wszystkich
procesów tej klasy i dla każdego $t\in[0,T]$.
\end{twierdzenie}
\begin{proof}
Gęstość procesów elementarnych otrzymujemy następująco.
Wystarczy najpierw rozważyć funkcje ograniczone, wspierane
na zbiorach skończonej miary $\Prob\otimes\dd t\otimes\nu$;
obcięcia wartości i lokalizacja w przestrzeni $E$ pozwalają
przejść do tego przypadku. Prostokąty przewidywalne
$F_s\times(s,t]\times A$ generują odpowiednie $\sigma$-ciało,
a więc ich liniowe kombinacje są gęste w $L^2$ z
argumentu klasy monotonicznej. Dla dowolnej aproksymacji
$H^{(n)}$ izometria daje zbieżność w $L^2$ zmiennych
$I_T(H^{(n)})$. Nierówność maksymalna Dooba zapewnia
\[
\E\sup_{t\le T}|I_t(H^{(n)})-I_t(H^{(m)})|^2
\le4\,\E\int_0^T\int_E|H^{(n)}-H^{(m)}|^2
\,\dd\nu\,\dd t.
\]
Wybieramy podciąg z sumowalnymi kwadratami różnic
w normie supremum, uzyskując przez Borela--Cantellego
jednostajną granicę prawie na pewno. Ponieważ jednostajna
granica funkcji càdlàg jest càdlàg, otrzymujemy taką
wersję $I(H)$. Jednoznaczność granicy w normie $L^2$
i nierówność maksymalna zapewniają niezależność od ciągu
aproksymującego i nierozróżnialność. Martyngałowość
wynika z granic w $L^2$, a izometria z ciągłości iloczynu
skalarnego w $L^2$. Zbieżność ucp zachodzi dla całej
aproksymacji, nie tylko podciągu.
\end{proof}

\begin{wniosek}[Nawias przewidywalny i kowariacja]
\label{cor:prm-angle}
Dla $H,K$ jak w Twierdzeniu \ref{th:prm-extend} zachodzi
\[
\langle I(H),I(K)\rangle_t
=\int_0^t\int_EH_s(z)K_s(z)\,\dd\nu(z)\,\dd s.
\]
W szczególności proces $I_t(H)^2-\int_0^t\int_EH_s(z)^2
\,\dd\nu(z)\,\dd s$ jest martyngałem.
\end{wniosek}
\begin{proof}
Dla elementarnego $H$ obliczamy przyrosty kwadratu.
Po warunkowaniu względem przeszłości przyrost kwadratu
części martyngałowej ma średnią
$\int_{t_j}^{t_{j+1}}\int_EH_s(z)^2\,\dd\nu\,\dd s$;
wyrazy mieszane mają średnią zero. Dla ogólnego $H$
wybieramy aproksymację w normie $L^2$. Izometria i
nierówność Cauchy'ego--Schwarza dają zbieżność
$I_t(H^{(n)})^2\to I_t(H)^2$ w $L^1$, a
$\int_0^t\int(H^{(n)})^2\dd\nu\dd s\to
\int_0^t\int H^2\dd\nu\dd s$ w $L^1$.
Warunkowa definicja martyngału przechodzi więc do granicy.
Tożsamość z kowariacją dostajemy przez polaryzację.
\end{proof}

\section{Skoki całki, wariacja kwadratowa i interpretacja}

\begin{twierdzenie}[Wariacja kwadratowa całki skompensowanej]
\label{th:prm-bracket}
Dla przewidywalnego $H$ spełniającego założenia
Twierdzenia \ref{th:prm-extend} i procesu $M=I(H)$ zachodzi
\[
[M]_t=\int_{(0,t]\times E}H_s(z)^2\,N(\dd s,\dd z),
\qquad
\langle M\rangle_t=\int_0^t\int_EH_s(z)^2\,\dd\nu(z)\,\dd s.
\tag{11.3}\label{eq:prm-brackets}
\]
W szczególności $\Delta M_t=\int_EH_t(z)N(\{t\},\dd z)$
dla wszystkich $t$ poza jednym zbiorem zerowym.
\end{twierdzenie}
\begin{proof}
Najpierw ustalmy $t\le T$ oraz deterministyczny podział
$\pi=\{0=t_0<\cdots<t_m=t\}$; piszemy
$Q_\pi(Y)=\sum_{j=1}^m(Y_{t_j}-Y_{t_{j-1}})^2$.
Dla procesu elementarnego $H^{(n)}$, ograniczonego i
wspieranego na zbiorze o skończonej intensywności, całka
$I(H^{(n)})$ ma na $[0,t]$ skończenie wiele skoków, a między
nimi zmienia się w sposób absolutnie ciągły. Z rozdzielenia
przedziałów zawierających skoki i ciągłych odcinków trajektorii
wynika, że dla każdego ciągu podziałów o średnicy dążącej do zera
\[
 Q_\pi(I(H^{(n)}))\longrightarrow
 \int_{(0,t]\times E}(H_s^{(n)}(z))^2\,N(\dd s,\dd z)
 \quad\text{prawie na pewno}.
\]
To jest własność wariacji kwadratowej funkcji o skończonej
wariacji złożonej ze skończenie wielu skoków, a nie tylko
stwierdzenie o zbieżności procesów w supremum.

Niech teraz $H^{(n)}\to H$ w
$L^2(\Prob\otimes\dd s\otimes\nu)$ i oznaczmy
$M^{(n)}=I(H^{(n)})$, $M=I(H)$. Dla dowolnego martyngału
kwadratowo całkowalnego $D$ o $D_0=0$ ortogonalność
przyrostów daje
$\E Q_\pi(D)=\E D_t^2$. Stosując
$|a^2-b^2|\le |a-b|\,|a+b|$ i nierówność
Cauchy'ego--Schwarza do sum po podziale, otrzymujemy
\[
\begin{split}
\E|Q_\pi(M)-Q_\pi(M^{(n)})|
&\le (\E Q_\pi(M-M^{(n)}))^{1/2}
      (\E Q_\pi(M+M^{(n)}))^{1/2}\\
&\le \norm{H-H^{(n)}}_2
      (\norm H_2+\norm{H^{(n)}}_2).
\end{split}
\]
Oszacowanie jest \emph{jednostajne względem podziału}.
Ponadto formuła kompensacji i ponowne zastosowanie
Cauchy'ego--Schwarza dają
\[
\E\left|\int ((H^{(n)})^2-H^2)\,\dd N\right|
\le\norm{H-H^{(n)}}_2
  (\norm H_2+\norm{H^{(n)}}_2)\longrightarrow0.
\]
Przyjmując najpierw $n$ duże, a następnie dzieląc coraz
drobniej przedział $[0,t]$, dowodzimy więc, że
$Q_\pi(M)\to\int H^2\dd N$ w prawdopodobieństwie.
Ta granica jest wariacją kwadratową $[M]_t$.

W przypadku procesów elementarnych, rozwijając kwadrat
skompensowanych przyrostów, otrzymujemy
$M_t^2-\int_{(0,t]\times E}H_s(z)^2N(\dd s,\dd z)$
jako martyngał. Dla ogólnego $H$ przejście do granicy jest
uprawnione w $L^1$: izometria kontroluje kwadraty $M_t^{(n)}$,
a poprzednie oszacowanie kontroluje całki względem $N$.
Nawias przewidywalny wynika z Wniosku~\ref{cor:prm-angle}.
Tożsamość dla skoków uzyskujemy najpierw dla procesów
elementarnych, a następnie wzdłuż podciągu jednostajnie
zbieżnego prawie na pewno oraz przez aproksymację całki
względem $N$ na zbiorach o skończonej intensywności.
\end{proof}

\begin{uwaga}[Dlaczego nawiasy są różne?]
Dla $M_t=N_t-\lambda t$, gdzie $N$ jest procesem
Poissona, mamy $[M]_t=N_t$, ale
$\langle M\rangle_t=\lambda t$. Pierwszy proces
jest losowym \emph{zrealizowanym} sumowaniem kwadratów
skoków, drugi jego przewidywalną średnią.
Ten drugi pojawia się w izometrii $L^2$.
\end{uwaga}

\begin{zadanie}[Skoki dodatnie i ujemne, wyprowadzenie dwóch nawiasów]
Niech $N^+,N^-$ będą niezależnymi procesami Poissona z
intensywnościami $\lambda_+,\lambda_->0$. Dla
$M_t=(N_t^+-\lambda_+t)-(N_t^--\lambda_-t)$ obliczyć
$[M]_t$ i $\langle M\rangle_t$.

\emph{Rozwiązanie.} Bierzemy $E=\{-1,1\}$,
$\nu=\lambda_+\delta_1+\lambda_-\delta_{-1}$,
$H(t,z)=z$. Przy prawdopodobieństwie jeden zliczenia
$N^+,N^-$ nie skaczą w tej samej chwili. Zatem
$[M]_t=\sum_{s\le t}(\Delta M_s)^2=N_t^++N_t^-$ oraz
$\langle M\rangle_t=(\lambda_++\lambda_-)t$.
Tożsamości weryfikuje też wariancja
$\E M_t^2=t(\lambda_++\lambda_-)$.
\end{zadanie}

\section{Lokalizacja i skoki o nieskończonej aktywności}

\begin{twierdzenie}[Lokalna całkowalność kwadratowa]
\label{th:prm-local}
Jeżeli $H$ jest przewidywalna i
$\int_0^T\int_EH_t(z)^2\,\dd\nu(z)\,\dd t<\infty$
prawie na pewno, to istnieje càdlàg lokalny martyngał
$\int_0^t\int_EH_s(z)\,\widetilde N(\dd s,\dd z)$,
który po zatrzymaniu w odpowiednim ciągu czasów
jest całką z Twierdzenia \ref{th:prm-extend}.
\end{twierdzenie}
\begin{proof}
Niech $A_t=\int_0^t\int_EH_s(z)^2\,\dd\nu(z)\,\dd s$.
Jest procesem ciągłym, rosnącym i adaptowanym.
Definiujemy $\tau_n=\inf\{t:A_t>n\}\wedge T$.
Ciągłość $A$ daje $A_{\tau_n}\le n$, więc proces
$H\ind{\{t\le\tau_n\}}$ należy do przestrzeni $L^2$
z normą nie większą od $n$. Własność zatrzymywania dla
całki skompensowanej wynika najpierw z definicji dla
procesów elementarnych, a następnie przez aproksymację
w normie $L^2$. Całki na kolejnych poziomach lokalizacji
są zgodne do chwili $\tau_n$, więc sklejają się w jeden
proces, określony z dokładnością do nierozróżnialności.
Ponieważ $A_T<\infty$ prawie na pewno, dla każdej
trajektorii $\tau_n=T$ od pewnego $n$.
\end{proof}

\begin{przyklad}[Miara stabilna i dlaczego kompensacja jest potrzebna]
Niech w jednym wymiarze $\dd\nu(z)=c|z|^{-1-\alpha}\dd z$
na $\{0<|z|\le1\}$, gdzie $0<\alpha<2$ i $c>0$.
Wtedy $\int_{|z|\le1}z^2\,\dd\nu(z)=2c/(2-\alpha)$,
więc $M_t=\int_0^t\int_{|z|\le1}z\,\widetilde N(\dd s,\dd z)$
istnieje w $L^2$. Jeśli $\alpha\ge1$, to
$\int_{|z|\le1}|z|\,\dd\nu(z)=\infty$; nie możemy
zastąpić całki skompensowanej zwykłą sumą bezwzględną
skoków. Mimo to jej wariacja kwadratowa jest skończona
prawie na pewno na $[0,T]$, bo
$\E[M]_T=T\,2c/(2-\alpha)<\infty$.
\end{przyklad}

\begin{zadanie}[Czekanie na pierwszy skok i losowa funkcja podcałkowa]
Niech $N_t$ będzie jednorodnym procesem Poissona o intensywności
$\lambda>0$, a $\tau=\inf\{t>0:N_t\ge1\}$ jego pierwszą chwilą
skoku. Zdefiniujmy $H_t=\ind{\{N_{t-}=0\}}$.
Wyznaczyć $M_t=\int_0^tH_s\,\dd(N_s-\lambda s)$, jego
średnią i drugi moment.

\emph{Rozwiązanie.} Lewostronnie ciągły, adaptowany proces
$H$ jest przewidywalny. Do pierwszego skoku jest równy $1$,
a po pierwszym skoku jest równy $0$. W całce względem
$\dd N$ występuje więc dokładnie jeden skok wielkości $1$,
jeśli $\tau\le t$. W całce względem czasu otrzymujemy
długość odcinka przed pierwszym skokiem. Stąd
\[
M_t=\ind{\{\tau\le t\}}-\lambda(t\wedge\tau).
\]
Ponieważ $\Prob(\tau>s)=e^{-\lambda s}$,
$\E(t\wedge\tau)=\int_0^t\Prob(\tau>s)\,\dd s
=(1-e^{-\lambda t})/\lambda$.
Wynika stąd $\E M_t=0$. Z izometrii Poissona,
\[
\E M_t^2=\lambda\E\int_0^tH_s^2\,\dd s
=\lambda\E(t\wedge\tau)=1-e^{-\lambda t}.
\]
To przykład całki z \emph{losowym} współczynnikiem,
który można obliczyć bez znajomości rozkładu całego procesu
całkowego. Próba użycia $\ind{\{N_t=0\}}$ jako współczynnika
na całym przedziale $(0,t]$ korzystałaby z informacji
przyszłej i nie spełniałaby tego argumentu.
\end{zadanie}

\section{Dekompozycja Lévy'ego--Itô: dokładne znaczenie formuły}

\begin{twierdzenie}[Reprezentacja konstruowanego procesu Lévy'ego]
\label{th:levy-ito-made}
Niech $b\in\R^d$, $Q$ będzie symetryczną macierzą
nieujemnie określoną, a $\nu$ miarą Lévy'ego.
Na pewnej przestrzeni probabilistycznej można zbudować
niezależne: $d$-wymiarowy proces Wienera $W$ i miarę
Poissona $N$ na $(0,\infty)\times(\R^d\setminus\{0\})$
o intensywności $\dd t\,\dd\nu(z)$, tak że proces
\[
X_t=bt+Q^{1/2}W_t
+\int_0^t\int_{\|z\|\le1}z\,\widetilde N(\dd s,\dd z)
+\int_0^t\int_{\|z\|>1}z\,N(\dd s,\dd z)
\tag{11.4}\label{eq:levy-ito-made}
\]
ma wykładnik Lévy'ego--Chinczyna odpowiadający trójce
$(b,Q,\nu)$. Całka dużych skoków jest sumą skończenie
wielu wyrazów na każdym przedziale zwartym, nawet jeśli
nie ma skończonej wartości oczekiwanej.
\end{twierdzenie}
\begin{proof}
Z Twierdzenia \ref{th:prm-existence} bierzemy miarę $N$;
z konstrukcji produktowej uniezależniamy od niej $W$.
Ponieważ $\int_{\|z\|\le1}\|z\|^2\,\dd\nu(z)<\infty$,
mała całka skompensowana istnieje na każdym horyzoncie,
a z \eqref{eq:prm-isometry} jest ograniczona w $L^2$.
Ponieważ $\nu(\{\|z\|>1\})<\infty$, zliczenia dużych
skoków są lokalnie skończone. Dla $\varepsilon>0$
zastąpmy małą całkę przez skompensowaną całkę po
$\varepsilon<\|z\|\le1$. Jest to złożony proces
Poissona pomniejszony o jego liniowy średni dryf.
Jego wykładnik, z rachunku w Rozdziale \ref{ch:levy},
wynosi $t\int_{\varepsilon<\|z\|\le1}
(e^{i\langle u,z\rangle}-1-i\langle u,z\rangle)\,\dd\nu(z)$.
Izometria i nierówność Dooba zapewniają zbieżność ucp
małych całek przy $\varepsilon\downarrow0$.
Ponieważ $|e^{iv}-1-iv|\le v^2/2$, po przejściu
do granicy otrzymujemy małoskokowy składnik wykładnika.
Dodając niezależne części: $bt$, $Q^{1/2}W$ i duże skoki,
dostajemy dokładnie formułę \eqref{eq:LK}.
Własności Lévy'ego wynikają z niezależnych stacjonarnych
przyrostów składników i prawostronnej ciągłości ich
trajektorii. Uzupełnia to konstrukcyjny kierunek
Twierdzenia \ref{th:LK}; kierunek klasyfikacyjny pozostaje
przyjętym twierdzeniem pomocniczym.
\end{proof}

\begin{zadanie}[Poisson i wykładnik charakterystyczny z całki]
Dla $\nu=\lambda\delta_a$ na $\R\setminus\{0\}$,
$\lambda>0$, $a\ne0$, obliczyć
$Y_t=\int_0^t\int z\,N(\dd s,\dd z)$ oraz jego wykładnik.

\emph{Rozwiązanie.} Jedynym dopuszczalnym rozmiarem skoku
jest $a$, a liczba skoków do chwili $t$ ma rozkład
Poissona o parametrze $\lambda t$. Zatem $Y_t=aN_t$ i
$\E e^{iuY_t}=\exp(\lambda t(e^{iua}-1))$.
Jeżeli $|a|\le1$, to w trójce z obcięciem $z\ind{\{|z|\le1\}}$
parametr $b=\lambda a$; jeśli $|a|>1$, mamy $b=0$.
To dlatego ta sama suma skoków ma inny zapis współczynnika
$b$ przy innej funkcji obcięcia.
\end{zadanie}

\begin{center}
\begin{tikzpicture}[>=Latex,scale=.9]
\draw[->] (0,0)--(9,0) node[right] {$t$};
\draw[->] (0,-1.3)--(0,3.2) node[above] {$M_t$};
\draw[blue!65!black,very thick]
(0,0)--(1.4,-.27)--(1.4,1.08)--(3,-.05)--(3,1.07)--
(4.6,.75)--(4.6,2.14)--(6.2,1.82)--(6.2,1.12)--(8.5,.72);
\draw[orange!80!black,thick,dashed] (0,0)--(8.4,-1.0);
\node[anchor=west,scale=.85,blue!65!black] at (5.6,2.75)
{skoki kompensowane};
\node[anchor=west,scale=.82,orange!80!black] at (4.9,-.92)
{przewidywalny trend};
\end{tikzpicture}
\par\smallskip\small Schemat procesu skokowego: odjęcie kompensatora
zmienia zachowanie między skokami, ale nie ich rozmiary.
Rysunek jest poglądowy, nie stanowi symulacji.
\end{center}

\section{Granica teorii i przygotowanie do półmartyngałów}
Zdefiniowaliśmy całkę względem skompensowanej miary
Poissona w klasie kwadratowo całkowalnej i przez lokalizację.
Udało się również uzasadnić budowę procesu z zadanej trójki
Lévy'ego bez ukrytego założenia o całkowalności dużych
skoków. Ogólne teoretyczne pojęcie \emph{kompensatora}
losowej miary skoków, której intensywność może zależeć
od przeszłości, wymaga teorii procesów przewidywalnych
i półmartyngałów; do niego wrócimy w rozdziale następnym.

% ===========================================================================
% ROZDZIAL 12 -- Polmartyngaly i wzor Ito ze skokami
% ===========================================================================
